% Part: lambda-calculus
% Chapter: introduction
% Section: lambda-computable

\documentclass[../../../include/open-logic-section]{subfiles}

\begin{document}

\olfileid{lam}{int}{cmp}
\olsection{\usetoken{S}{lambda definable} فنکشن قابلِ حساب نیں}

\begin{thm}
\ollabel{thm:lambda-computable}
جے اک جزوی فنکشن $f$ کسے لامبڈا ٹرم نال !!{lambda defined}
ہووے، تاں اوہ قابلِ حساب اے۔
\end{thm}

\begin{proof}
فرض کرو فنکشن~$f$ لامبڈا ٹرم~$X$ نال !!{lambda defined}
اے۔ آؤ $f$ دا حساب کرن دا اک غیر رسمی طریقہ بیان کریے۔
انپٹ $m_0$، \dots،~$m_{n-1}$ ملن اُتے ٹرم
$X \num{m_0} \ldots \num{m_{n-1}}$ لکھ لو۔اک درخت بناؤ:
پہلاں اصل ٹرم دیاں ساریاں اک قدم دیاں تقلیلاں لکھو۔ اوہناں دے
ہِٹھاں اوہناں وچوں ہر اک دیاں ساریاں اک قدم دیاں تقلیلاں لکھو
(یعنی اصل ٹرم دیاں دو-قدم تقلیلاں)؛فیر ایہ کم جاری رکھو۔
جے کدے کوئی چرچ عدد لبھ پئے تاں اوہی جواب دے طور اُتے واپس کرو۔
جے کدے نہ لبھے تاں تلاش جاری رہندی اے تے فنکشن اوس انپٹ اُتے متعین نہیں۔

چرچ دے دعوے دی ورتوں نال آکھ سکدے آں کہ ایہ فنکشن قابلِ
حساب اے۔قضیہ ثابت کرن دا بہتر طریقہ ایہ ہووے گا کہ ایس
تلاش دی بازگشتی وضاحت دتی جائے۔مثلاً،بنیادی بازگشتی فنکشناں
تے تعلقاں دا اک سلسلہ تعریف کیتا جا سکدا اے:
``$\fn{IsASubterm}$''، ``$\fn{Substitute}$''،
``$\fn{ReducesToInOneStep}$''، ``$\fn{ReductionSequence}$''،
``$\fn{Numeral}$'' وغیرہ۔$f(m_0, \dots, m_{n-1})$
دا حساب کرن دا جزوی بازگشتی طریقہ فیر
$X \num{m_0} \dots \num{m_{n-1}}$ توں شروع ہو کے
کسے چرچ عدد اُتے ختم ہون والی اک-قدم تقلیلاں دی لڑی
لبھدا اے،تے اوس چرچ عدد دے مطابق عدد واپس کردا اے۔
تفصیل لمبی تے اکتاؤنی اے،پر باقی کم معمول دے مطابق اے۔
\end{proof}

\end{document}
